Another question in Darvas--Xia--Zhang concerns the raw filtration test
curve for a big line bundle. Immediately before their Definition~4.4,
they formulate the equality
\(u_\tau^{\mathcal F}=P[u_\tau^{\mathcal F}]\) as a conjecture.
Under the hypotheses of their Definition~4.1, every strict-interior
parameter \(\tau<\tau_L(\mathcal F)\) has positive top nonpluripolar
mass. Their Lemma~4.5 already proves
\(P[u_\tau^{\mathcal F}]=P[u_\tau^{\mathcal F}]_{\mathcal I}\);
the additional question is equality with the raw function, not equality
of these two envelopes
\cite[Definition~4.3, the paragraph before Definition~4.4,
Lemma~4.5, and the proof of Proposition~4.6]{DXZ23}.
For a later approximation theorem for graded-series equilibrium metrics,
see \cite[Theorem~4.3]{Finski26Kodaira}.

A reconstruction question involving zero-top-mass data must be formulated
separately. It is not a second zero-mass interval strictly before
\(\tau_L(\mathcal F)\) in this construction, and arbitrary zero-mass
reconstruction must not be identified with a terminal-left-limit
observation. Such a question must specify the input data, the ambient
class, the normalized target, and the equivalence relation: these choices
are not supplied by the preceding filtration conjecture.
